\documentclass[11pt,letterpaper,hidelinks]{book}
\usepackage{setspace}
\usepackage{fancyhdr}
\usepackage{ragged2e}
\usepackage[utf8]{inputenc}
\usepackage[T1]{fontenc}
\usepackage{amsmath}
\usepackage{amsfonts,mathrsfs}
\usepackage{amssymb}
\usepackage{amsthm}
\usepackage{setspace}
\usepackage{enumitem}
\usepackage{comment}
\usepackage[spanish, es-tabla, es-nodecimaldot]{babel}
\usepackage{graphicx,color}
\usepackage{epsfig} 
\usepackage{hyperref}
\usepackage{float}
\usepackage{listings}
\usepackage{multirow}
\usepackage{caption}
\usepackage{subcaption}
\usepackage[numbered,framed]{matlab-prettifier}
\lstMakeShortInline"
\lstset{
  style              = Matlab-editor,
  escapechar         = ",
  mlshowsectionrules = true,
}

\usepackage[top=3cm,bottom=2cm,left=2cm,right=2cm]{geometry}
%\usepackage[letterpaper, total={6in, 8in}]{geometry}
%\pagestyle{headings}

\setlength{\parindent}{0pt}

\hypersetup{colorlinks, linkcolor=blue, citecolor=blue,urlcolor=blue,plainpages=false}

\newtheorem{theorem}{Teorema}[section]
\newtheorem{lemma}[theorem]{Lemma}

\renewcommand\qedsymbol{$\blacksquare$}

\theoremstyle{definition}
\newtheorem{problem}{Problema}


\newtheorem{affir}{Afirmación}
\newtheorem{defi}{Defición}
\newtheorem{corollary}[theorem]{Corollary}
\newtheorem{assumption}[theorem]{Assumption}
\newtheorem{remark}[theorem]{Remark}
\newtheorem{proposition}[theorem]{Proposition}
\newtheorem{algorithm}[theorem]{Algorithm}
\newtheorem{experiment}[theorem]{Experiment}
\usepackage{todonotes}
\usepackage{xcolor}
\usepackage{cancel}
\usepackage{empheq}
\newcommand{\pvr}[1]{{\color{blue}{#1}}}
\DeclareMathOperator*{\argmin}{arg\,min}
\usepackage{authblk}
\newcommand{\bpsi}{\boldsymbol\psi}
\newcommand{\mover}[1]{{\color{red}{#1}}}
\newcommand{\keko}[1]{{\color{magenta}{#1}}}

\newcommand{\liliana}[2]{{\color{violet}{#1}}}

\newcommand{\Cnueve}{D_1}
\newcommand{\Cdiez}{D_2}
\newcommand{\compD}{\Omega_h}
\newcommand{\vertiii}[1]{{\left\vert\kern-0.25ex\left\vert\kern-0.25ex\left\vert #1 
    \right\vert\kern-0.25ex\right\vert\kern-0.25ex\right\vert}}

\newcommand\revise[1]{\textcolor{blue}{#1}}
%\newcommand\comment[1]{\textcolor{red}{#1}}
\newcommand\commentM[1]{\textcolor{magenta}{#1}}

\newcommand\red[1]{\textcolor{red}{#1}}


\newcommand\Source{S(f,g)}
\newcommand\Sourceone{S(f,g,1)}


\definecolor{mygreen}{rgb}{0.0,0.5,0.0}


\usepackage{mathrsfs}
%% CODIGO MATLAB
\usepackage{listings}
\lstset{
    language=Matlab,
    basicstyle=\ttfamily\small,
    keywordstyle=\color{blue},
    commentstyle=\color{green!50!black},
    stringstyle=\color{red},
    showstringspaces=false,
    tabsize=4,
    breaklines=true,
    breakatwhitespace=true
}

\newcommand{\triple}[1]{|\!|\!|{#1}|\!|\!|}
%% Miscellaneous
\renewcommand{\Re}{\operatorname{Re}}
\renewcommand{\Im}{\operatorname{Im}}

\newcommand{\eq}{:=}

%% Differential operator
\newcommand{\pd}[2]{\frac{\partial #1}{\partial #2}}

\newcommand{\grad}{\boldsymbol \nabla}
\renewcommand{\div}{\grad \cdot}
\newcommand{\curl}{\grad \times}

\newcommand{\ccurl}{\boldsymbol{\operatorname{curl}}}
\newcommand{\ddiv}{\operatorname{div}}

%% DG operators
\newcommand{\jmp}[1]{\,[\![#1]\!]}
\newcommand{\avg}[1]{\{\!\!\{#1\}\!\!\}}

%% Bold uppercases
\newcommand{\BA}{\boldsymbol A}
\newcommand{\BB}{\boldsymbol B}
\newcommand{\BC}{\boldsymbol C}
\newcommand{\BD}{\boldsymbol D}
\newcommand{\BE}{\boldsymbol E}
\newcommand{\BF}{\boldsymbol F}
\newcommand{\BG}{\boldsymbol G}
\newcommand{\BH}{\boldsymbol H}
\newcommand{\BI}{\boldsymbol I}
\newcommand{\BJ}{\boldsymbol J}
\newcommand{\BK}{\boldsymbol K}
\newcommand{\BL}{\boldsymbol L}
\newcommand{\BM}{\boldsymbol M}
\newcommand{\BN}{\boldsymbol N}
\newcommand{\BO}{\boldsymbol O}
\newcommand{\BP}{\boldsymbol P}
\newcommand{\BQ}{\boldsymbol Q}
\newcommand{\BR}{\boldsymbol R}
\newcommand{\BS}{\boldsymbol S}
\newcommand{\BT}{\boldsymbol T}
\newcommand{\BU}{\boldsymbol U}
\newcommand{\BV}{\boldsymbol V}
\newcommand{\BW}{\boldsymbol W}
\newcommand{\BX}{\boldsymbol X}
\newcommand{\BY}{\boldsymbol Y}
\newcommand{\BZ}{\boldsymbol Z}

%% Bold lowercases
\newcommand{\ba}{\boldsymbol a}
\newcommand{\bb}{\boldsymbol b}
\newcommand{\bc}{\boldsymbol c}
\newcommand{\bd}{\boldsymbol d}
\newcommand{\be}{\boldsymbol e}
\newcommand{\bff}{\boldsymbol f}
\newcommand{\bg}{\boldsymbol g}
\newcommand{\bh}{\boldsymbol h}
\newcommand{\bi}{\boldsymbol i}
\newcommand{\bj}{\boldsymbol j}
\newcommand{\bk}{\boldsymbol k}
\newcommand{\bl}{\boldsymbol l}
%\newcommand{\bm}{\boldsymbol m}
\newcommand{\bn}{\boldsymbol n}
\newcommand{\bo}{\boldsymbol o}
\newcommand{\bp}{\boldsymbol p}
\newcommand{\bq}{\boldsymbol q}
\newcommand{\br}{\boldsymbol r}
\newcommand{\bs}{\boldsymbol s}
\newcommand{\bt}{\boldsymbol t}
\newcommand{\bu}{\boldsymbol u}
\newcommand{\bv}{\boldsymbol v}
\newcommand{\bw}{\boldsymbol w}
\newcommand{\bx}{\boldsymbol x}
\newcommand{\by}{\boldsymbol y}
\newcommand{\bz}{\boldsymbol z}

%% Curly uppercases
\newcommand{\CA}{\mathcal A}
\newcommand{\CB}{\mathcal B}
\newcommand{\CC}{\mathcal C}
\newcommand{\CD}{\mathcal D}
\newcommand{\CE}{\mathcal E}
\newcommand{\CF}{\mathcal F}
\newcommand{\CG}{\mathcal G}
\newcommand{\CH}{\mathcal H}
\newcommand{\CI}{\mathcal I}
\newcommand{\CJ}{\mathcal J}
\newcommand{\CK}{\mathcal K}
\newcommand{\CL}{\mathcal L}
\newcommand{\CM}{\mathcal M}
\newcommand{\CN}{\mathcal N}
\newcommand{\CO}{\mathcal O}
\newcommand{\CP}{\mathcal P}
\newcommand{\CQ}{\mathcal Q}
\newcommand{\CR}{\mathcal R}
\newcommand{\CS}{\mathcal S}
\newcommand{\CT}{\mathcal T}
\newcommand{\CU}{\mathcal U}
\newcommand{\CV}{\mathcal V}
\newcommand{\CW}{\mathcal W}
\newcommand{\CX}{\mathcal X}
\newcommand{\CY}{\mathcal Y}
\newcommand{\CZ}{\mathcal Z}

%% Caligraphic uppercases
\newcommand{\LA}{\mathscr A}
\newcommand{\LB}{\mathscr B}
\newcommand{\LC}{\mathscr C}
\newcommand{\LD}{\mathscr D}
\newcommand{\LE}{\mathscr E}
\newcommand{\LF}{\mathscr F}
\newcommand{\LG}{\mathscr G}
\newcommand{\LH}{\mathscr H}
\newcommand{\LI}{\mathscr I}
\newcommand{\LJ}{\mathscr J}
\newcommand{\LK}{\mathscr K}
\newcommand{\LL}{\mathscr L}
\newcommand{\LM}{\mathscr M}
\newcommand{\LN}{\mathscr N}
\newcommand{\LO}{\mathscr O}
\newcommand{\LP}{\mathscr P}
\newcommand{\LQ}{\mathscr Q}
\newcommand{\LR}{\mathscr R}
\newcommand{\LS}{\mathscr S}
\newcommand{\LT}{\mathscr T}
\newcommand{\LU}{\mathscr U}
\newcommand{\LV}{\mathscr V}
\newcommand{\LW}{\mathscr W}
\newcommand{\LX}{\mathscr X}
\newcommand{\LY}{\mathscr Y}
\newcommand{\LZ}{\mathscr Z}

%% Caligraphic uppercases
\newcommand{\TA}{\textup A}
\newcommand{\TB}{\textup B}
\newcommand{\TC}{\textup C}
\newcommand{\TD}{\textup D}
\newcommand{\TE}{\textup E}
\newcommand{\TF}{\textup F}
\newcommand{\TG}{\textup G}
%\newcommand{\TH}{\textup H}
\newcommand{\TI}{\textup I}
\newcommand{\TJ}{\textup J}
\newcommand{\TK}{\textup K}
\newcommand{\TL}{\textup L}
\newcommand{\TM}{\textup M}
\newcommand{\TN}{\textup N}
\newcommand{\TO}{\textup O}
\newcommand{\TP}{\textup P}
\newcommand{\TQ}{\textup Q}
\newcommand{\TR}{\textup R}
\newcommand{\TS}{\textup S}
\newcommand{\TT}{\textup T}
\newcommand{\TU}{\textup U}
\newcommand{\TV}{\textup V}
\newcommand{\TW}{\textup W}
\newcommand{\TX}{\textup X}
\newcommand{\TY}{\textup Y}
\newcommand{\TZ}{\textup Z}

%% Bold curly uppercases
\newcommand{\BCA}{\boldsymbol{\CA}}
\newcommand{\BCB}{\boldsymbol{\CB}}
\newcommand{\BCC}{\boldsymbol{\CC}}
\newcommand{\BCD}{\boldsymbol{\CD}}
\newcommand{\BCE}{\boldsymbol{\CE}}
\newcommand{\BCF}{\boldsymbol{\CF}}
\newcommand{\BCG}{\boldsymbol{\CG}}
\newcommand{\BCH}{\boldsymbol{\CH}}
\newcommand{\BCI}{\boldsymbol{\CI}}
\newcommand{\BCJ}{\boldsymbol{\CJ}}
\newcommand{\BCK}{\boldsymbol{\CK}}
\newcommand{\BCL}{\boldsymbol{\CL}}
\newcommand{\BCM}{\boldsymbol{\CM}}
\newcommand{\BCN}{\boldsymbol{\CN}}
\newcommand{\BCO}{\boldsymbol{\CO}}
\newcommand{\BCP}{\boldsymbol{\CP}}
\newcommand{\BCQ}{\boldsymbol{\CQ}}
\newcommand{\BCR}{\boldsymbol{\CR}}
\newcommand{\BCS}{\boldsymbol{\CS}}
\newcommand{\BCT}{\boldsymbol{\CT}}
\newcommand{\BCU}{\boldsymbol{\CU}}
\newcommand{\BCV}{\boldsymbol{\CV}}
\newcommand{\BCW}{\boldsymbol{\CW}}
\newcommand{\BCX}{\boldsymbol{\CX}}
\newcommand{\BCY}{\boldsymbol{\CY}}
\newcommand{\BCZ}{\boldsymbol{\CZ}}

%% Bold caligraphic uppercases
\newcommand{\BLA}{\pmb{\LA}}
\newcommand{\BLB}{\pmb{\LB}}
\newcommand{\BLC}{\pmb{\LC}}
\newcommand{\BLD}{\pmb{\LD}}
\newcommand{\BLE}{\pmb{\LE}}
\newcommand{\BLF}{\pmb{\LF}}
\newcommand{\BLG}{\pmb{\LG}}
\newcommand{\BLH}{\pmb{\LH}}
\newcommand{\BLI}{\pmb{\LI}}
\newcommand{\BLJ}{\pmb{\LJ}}
\newcommand{\BLK}{\pmb{\LK}}
\newcommand{\BLL}{\pmb{\LL}}
\newcommand{\BLM}{\pmb{\LM}}
\newcommand{\BLN}{\pmb{\LN}}
\newcommand{\BLO}{\pmb{\LO}}
\newcommand{\BLP}{\pmb{\LP}}
\newcommand{\BLQ}{\pmb{\LQ}}
\newcommand{\BLR}{\pmb{\LR}}
\newcommand{\BLS}{\pmb{\LS}}
\newcommand{\BLT}{\pmb{\LT}}
\newcommand{\BLU}{\pmb{\LU}}
\newcommand{\BLV}{\pmb{\LV}}
\newcommand{\BLW}{\pmb{\LW}}
\newcommand{\BLX}{\pmb{\LX}}
\newcommand{\BLY}{\pmb{\LY}}
\newcommand{\BLZ}{\pmb{\LZ}}

%% Bold text-up
\newcommand{\BTA}{\pmb{\TA}}
\newcommand{\BTB}{\pmb{\TB}}
\newcommand{\BTC}{\pmb{\TC}}
\newcommand{\BTD}{\pmb{\TD}}
\newcommand{\BTE}{\pmb{\TE}}
\newcommand{\BTF}{\pmb{\TF}}
\newcommand{\BTG}{\pmb{\TG}}
%\newcommand{\BTH}{\pmb{\TH}}
\newcommand{\BTI}{\pmb{\TI}}
\newcommand{\BTJ}{\pmb{\TJ}}
\newcommand{\BTK}{\pmb{\TK}}
\newcommand{\BTL}{\pmb{\TL}}
\newcommand{\BTM}{\pmb{\TM}}
\newcommand{\BTN}{\pmb{\TN}}
\newcommand{\BTO}{\pmb{\TO}}
\newcommand{\BTP}{\pmb{\TP}}
\newcommand{\BTQ}{\pmb{\TQ}}
\newcommand{\BTR}{\pmb{\TR}}
\newcommand{\BTS}{\pmb{\TS}}
\newcommand{\BTT}{\pmb{\TT}}
\newcommand{\BTU}{\pmb{\TU}}
\newcommand{\BTV}{\pmb{\TV}}
\newcommand{\BTW}{\pmb{\TW}}
\newcommand{\BTX}{\pmb{\TX}}
\newcommand{\BTY}{\pmb{\TY}}
\newcommand{\BTZ}{\pmb{\TZ}}


\def\ds{\displaystyle}
\def\l{\left(}
\def\r{\right)}
\def\labs{\left|}
\def\rabs{\right|}
\providecommand{\norm}[1]{\left\| #1 \right\|}
\providecommand{\abs}[1]{\labs #1 \rabs}
\providecommand{\pa}[1]{\left( #1 \right)}
\providecommand{\cor}[1]{\left[ #1 \right]}
\providecommand{\lla}[1]{\left\{ #1 \right\}}
\providecommand{\inner}[2]{\l #1,#2 \r}
\providecommand{\pdual}[2]{\left\langle #1,#2 \right\rangle}

% conjuntos de numeros
\def\HH{{\mathcal H}}
\def\R{\mathbb{R}}
\def\LL{\mathcal{L}}
\def\KK{\mathbb{K}}
\def\K{\mathcal{K}}
\def\C{\mathbb{C}}
\def\G{\mathcal{G}}
\def\F{\mathcal{F}}
\def\P{{\mathcal P}}
\def\I{\mathbb{I}}
\def\N{\mathbb{N}}
\def\Q{\mathbb{Q}}
\def\Z{\mathbb{Z}}

\renewcommand{\theenumi}{(\alph{enumi})}
\newcommand{\balpha}{\boldsymbol{\alpha}}
\newcommand{\bbeta}{\boldsymbol{\beta}}

%%Simbolos tesis
\def\TT{\mathcal{T}}
\def\dTT{\partial \mathcal{T}}
\def\EE{\mathcal{E}}
\def\FF{\mathcal{F}}
\def\NN{\mathcal{N}}

\fancyhead[R]{\thepage}
\fancyfoot[C]{}

\begin{document}
\begin{titlepage}
\begin{center}
\includegraphics[scale=0.3]{img/esc.png} 
\end{center}
\newcommand{\HRule}{\rule{\linewidth}{0.5mm}} % Defines a new command for the horizontal lines, change thickness here
\center % Center everything on the page
%----------------------------------------------------------------------------------------
%	HEADING SECTIONS
%----------------------------------------------------------------------------------------
\textsc{\Large Universidad de Concepci\'on}\\ % Name of your university/college
\textsc{\large Facultad de Ciencias F\'isicas y Matem\'aticas}\\
\textsc{\large Departamento de Ingenier\'ia Matem\'atica}\\[2cm]
%----------------------------------------------------------------------------------------
%	TITLE SECTION
%----------------------------------------------------------------------------------------
\setstretch{1.5}
\begin{center}
\hrule
\vspace{0.5cm}
{\bf \Large A QUASI-PERIODIC HDG METHOD FOR THE HELMHOLTZ EQUATION IN HETEROGENEOUS DOMAINS}
\vspace{0.5cm}
\hrule
\end{center}
%----------------------------------------------------------------------------------------
%	AUTHOR SECTION
%----------------------------------------------------------------------------------------
\vspace{1cm}
\textsc{\large Autor}\\[0.3cm]
\Large{\bf Ignacio Sebastian Ruminot Aburto}\\
\vspace{2cm}
\setstretch{1}
Tesis presentada a la Facultad de Ciencias F\'isicas y Matem\'aticas de la
Universidad de Concepci\'on para optar al t\'itulo profesional de
Ingeniero Civil Matem\'atico\\[0.5cm]
\setstretch{1.5}
{Profesores Gu\'ia: Dr.~Manuel E.~Solano Dr.~Patrick~Vega.}\\[0.5cm]
\setstretch{1}
{\bf Septiembre de 2024,}
{\bf Concepci\'on, Chile.}
\end{titlepage}

\setstretch{1.5}
\clearpage
\setcounter{page}{1}
\thispagestyle{plain}
\hspace{1mm}
\vfill
\centering
{\copyright\hspace{1mm} 2024 Ignacio Sebastian Ruminot Aburto}\\[0.6cm]
\setstretch{1}
{Se autoriza la reproducci\'on total o parcial, con fines}
{acad\'emicos, por cualquier medio o procedimiento,}
{incluyendo la cita bibliogr\'afica del documento.}

\clearpage
\thispagestyle{plain}
\setstretch{1.5}
\center
{\bf \Large A QUASI-PERIODIC HDG METHOD FOR THE HELMHOLTZ EQUATION IN HETEROGENEOUS DOMAINS}
\vspace{1cm}
\justifying
\vspace{2.5cm}
{\bf \large COMISI\'ON EVALUADORA}
\vspace{1cm}
{Dr.~Manuel E.~Solano [Profesor gu\'ia]}
{CI$^2$MA y Departamento de Ingenier\'ia Matem\'atica, Universidad de Concepción, Chile.} 
\vspace{0.05cm}
{Dr.~Patrick Vega [Profesor gu\'ia]}
{Departamento de Matemática y Ciencia de la Computación, Universidad de Santiago de Chile, Chile.}
\vspace{0.2cm}
%{Dr.~Rommel Bustinza}
%{CI$^2$MA y Departamento de Ingenier\'ia Matem\'atica, Universidad de Concepción, Chile.}
\vspace{0.2cm}
%{Dr.~Paulo Zúñiga}
%{Department of Applied Mathematics, University of Waterloo, Canadá.}
\vspace{0.2cm}
{\it\bf FECHA DE DEFENSA: .}


\clearpage
\thispagestyle{plain}
\setstretch{1}
\center{\bf \LARGE Acknowledgements}
\phantomsection

\addcontentsline{toc}{chapter}{Acknowledgements}
\vspace{1cm}
\justifying
AQUÍ VAN LOS AGRADECIMIENTOS


% \clearpage
% \thispagestyle{plain}
% \justifying
% \begin{redarea}
% {\it Ackwnowledgements 2
% \vfill
% \begin{center}
% \includegraphics[scale=0.7]{images/logo_dim.png}
% \hspace{3cm}
% \includegraphics[scale=0.7]{images/logo_ci2ma.jpg} 

% \end{center}
% }
% \end{redarea}


\clearpage
\pagestyle{fancy}
\phantomsection
\tableofcontents

%  \clearpage
%  \addcontentsline{toc}{chapter}{List of figures} 
% \listoffigures

%  \clearpage
%  \addcontentsline{toc}{chapter}{List of tables} 
% \listoftables

\clearpage
\phantomsection
%\addcontentsline{toc}{chapter}{List of listings}
%\lstlistoflistings

\addcontentsline{toc}{chapter}{List of figures} 
\listoffigures


\clearpage

%**********************************************************************
%**********************************************************************
\thispagestyle{plain}
\center{\bf \LARGE Abstract}\\
\phantomsection
\addcontentsline{toc}{chapter}{Abstract}
\vspace{1cm}
\justifying

\clearpage

%**********************************************************************
%**********************************************************************
\thispagestyle{plain}
\center{\bf \LARGE Resumen}\\
\phantomsection
\addcontentsline{toc}{chapter}{Resumen}
\vspace{1cm}
\justifying


%**********************************************************
%**********************************************************
\clearpage
\pagestyle{fancy}
\setstretch{1.2}
\justifying
%\include{fuente/introduccion}

%\include{fuente/metodo_hdg}

%\include{fuente/analisis}

%\include{fuente/experimentos}

%\include{fuente/conclusion}


%\bibliography{fuente/bibliografia}
%\bibliographystyle{siam} 

hola

\end{document}